\documentclass[12pt,a4paper]{article}

\usepackage{courier}
\renewcommand{\familydefault}{\ttdefault}
\usepackage[margin=0.8in]{geometry}
\usepackage{enumitem}
\usepackage{setspace}
\usepackage{listings}
\usepackage{array}
\usepackage{fancyhdr}
\usepackage{tcolorbox}

\setlength{\parindent}{0pt}
\setlength{\parskip}{1.1em}
\setlength{\tabcolsep}{4pt}
\setlength{\extrarowheight}{3pt}
\pagestyle{plain}
\doublespacing

% Use - instead of bullet dot for lists, and set tight spacing
\setlist[itemize]{label=-, leftmargin=2em, nosep, itemsep=0.3em}

% Code listings style (white/light mode, courier font)
\lstset{
  backgroundcolor=\color{white},
  basicstyle=\ttfamily\small,
  breaklines=true,
  frame=single,
  framerule=0.5pt,
  rulecolor=\color{black},
  xleftmargin=10pt, xrightmargin=10pt,
  aboveskip=8pt, belowskip=8pt,
  showstringspaces=false,
  upquote=true,
  tabsize=4,
  columns=flexible
}

% Code block environment (verbatim matching Courier format)
\lstnewenvironment{codeblock}
  {\lstset{aboveskip=0pt,belowskip=0pt}}
  {}

% Table helpers
\newcolumntype{L}[1]{>{\raggedright\arraybackslash}p{#1}}
\newcolumntype{C}[1]{>{\centering\arraybackslash}p{#1}}
\newcommand{\thfmt}[1]{\textbf{#1}}
\newcommand{\tcfmt}[1]{#1}

\begin{document}
\raggedright

============================================================\\
\textbf{NULLSPLOIT TECHNICAL DOCUMENTATION}\\
\textbf{KERNEL FIREWALL \& SECURITY OPERATIONS CENTER}\\
\textbf{TECHNICAL ARCHITECTURE \& OPERATIONS MANUAL}\\
============================================================

\textbf{CLASSIFICATION:}\\
Internal Use Only

\textbf{DETAILS:}\\
Version: 1.0\\
Date: July 2026\\
Platform: Linux Netfilter/NFQUEUE, Docker\\
Status: Approved Documentation

\vspace{1.5em}
\textbf{1. EXECUTIVE SUMMARY}\\
------------------------------------------------------------\\
Nullsploit is a hybrid kernel-userspace network security stack designed to protect Apache web servers from common Layer 4 and Layer 7 network attacks. The system intercepts traffic using Linux's \textbf{Netfilter NFQUEUE} framework, inspects packets for known threat signatures, and forwards encrypted telemetry to a centralised SOC dashboard for real-time monitoring and historical analysis.

The stack has four tightly-integrated components: a userspace NFQUEUE firewall daemon (\texttt{fw\_nfq}), a telemetry agent (\texttt{fw\_agent}), a PGP-encrypted ingest receiver (\texttt{soc\_receiver}), and a real-time web dashboard. These are packaged as isolated Docker containers.

\textbf{1.1 Key Capabilities}\\
Below is a summary of the technical defense mechanisms and architectural features provided by the Nullsploit stack:

\begin{center}
\begin{tabular}{L{4.5cm} L{10cm}}
\hline
\thfmt{Capability} & \thfmt{Description} \\
\hline
PHP Shell Detection  & Pattern matching against RCE function signatures in URI and body \\
XSS Detection        & Case-insensitive scan for cross-site scripting vectors \\
SYN Flood Mitigation & Sliding-window rate limiting on TCP SYN packets per source IP \\
Slowloris Detection  & Timeout-based tracking of incomplete HTTP header connections \\
IP Blocklist         & CIDR-range lookups in a kernel-side hash table; hot-reloadable \\
PGP Encryption       & All agent-to-SOC telemetry encrypted with GPG public-key crypto \\
Live Dashboard       & SSE live event feed with Chart.js visualisations on port 8443 \\
Secure Auth          & bcrypt passwords, HMAC-SHA256 signed session tokens \\
\hline
\end{tabular}
\end{center}

\vspace{1.5em}
\textbf{2. SYSTEM ARCHITECTURE OVERVIEW}\\
------------------------------------------------------------\\
The operational flow begins at the packet interception layer and moves up through userspace analysis to secure telemetry transmission and graphical presentation.

\begin{center}
\begin{tabular}{L{3.5cm} L{3.5cm} L{7.5cm}}
\hline
\thfmt{Layer} & \thfmt{Component} & \thfmt{Role} \\
\hline
Network L4/L7 & iptables NFQUEUE & Redirect TCP/80 to queue 0 for userspace inspection \\
Userspace     & fw\_nfq Daemon   & Dequeue packets, inspect, issue NF\_ACCEPT or NF\_DROP \\
IPC           & UNIX Socket      & Transfer fw\_event structs from fw\_nfq to fw\_agent \\
Telemetry     & fw\_agent        & PGP-encrypt; send instantly (WARNING+) or batch (INFO) \\
Transport     & TCP Port 1113    & Encrypted channel to the SOC container \\
Ingest        & soc\_receiver    & epoll loop, decrypts PGP payloads, writes to SQLite \\
Storage       & SQLite           & Persistent event store at /var/lib/soc/soc.db \\
Presentation  & Web Dashboard    & REST API, SSE live feed, Chart.js on port 8443 \\
\hline
\end{tabular}
\end{center}

\textbf{2.1 Container Network Topology}\\
Both containers share the Docker bridge network. Service discovery uses container names as DNS hostnames. PGP keys are exchanged via a shared host volume at \texttt{./keys}.

\newpage
\textbf{2.2 Architecture Diagram}\\
The visual flow below details packet redirection and telemetry pathways:

\begin{codeblock}
  [ Attacker / Internet ]
          |  TCP:80
          v
  +-----------------------------------+
  |     APACHE SERVER CONTAINER       |
  |  iptables -> NFQUEUE (queue 0)    |
  |        fw_nfq daemon              |
  |     (inspect packet payloads)     |
  |              |  UNIX IPC          |
  |          fw_agent                 |
  |      (PGP encrypt & send)         |
  +-------------+---------------------+
                |   TCP:1113  (PGP)
                v
  +-----------------------------------+
  |      SOC SERVER CONTAINER         |
  |       soc_receiver daemon         |
  |   (epoll / decrypt / SQLite)      |
  |   SQLite DB  |  Web API :8443     |
  |         Nullsploit Dashboard      |
  +-----------------------------------+
\end{codeblock}

\vspace{1.5em}
\textbf{3. COMPONENT: KERNEL FIREWALL (\texttt{fw\_nfq})}\\
------------------------------------------------------------\\
The daemon processes incoming frames bound to NFQUEUE queue 0. The raw packet capture is achieved using the iptables rule configured at system initialization:

\begin{codeblock}
iptables -A INPUT -p tcp --dport 80 -j NFQUEUE --queue-num 0
\end{codeblock}

Each packet is evaluated sequentially against the following processing table:

\begin{center}
\begin{tabular}{C{1.2cm} L{5.2cm} L{8.1cm}}
\hline
\thfmt{Step} & \thfmt{Function} & \thfmt{Action on Detection} \\
\hline
1 & inspect\_\allowbreak ip\_\allowbreak blocklist() & Immediately issue NF\_DROP \\
2 & monitor\_\allowbreak tcp\_\allowbreak stats()    & NF\_DROP on SYN flood or Slowloris timeout \\
3 & inspect\_\allowbreak http\_\allowbreak payload() & NF\_DROP on PHP shell or XSS match \\
4 & (none)                   & Issue NF\_ACCEPT --- forward to Apache \\
\hline
\end{tabular}
\end{center}

\textbf{3.1 Telemetry Struct (\texttt{fw\_event})}\\
When a signature match or threshold violation is verified, the event parameters are packaged using the following binary structure and transmitted over a local UNIX socket at \texttt{/var/run/fw\_inspect.sock}:

\begin{codeblock}
struct fw_event {
    uint64_t  timestamp;
    uint32_t  src_ip;          /* Source IP (network byte order)        */
    uint32_t  dest_ip;
    uint16_t  src_port;
    uint16_t  dest_port;
    uint8_t   threat_type;     /* 1=PHP_SHELL 2=XSS 3=SYN_FLOOD 4=SLOWLORIS */
    uint8_t   severity;        /* 0=INFO  1=WARNING  2=CRITICAL         */
    char      payload_preview[128];
    char      details[256];
};
\end{codeblock}

\vspace{1.5em}
\textbf{4. COMPONENT: USERSPACE AGENT (\texttt{fw\_agent})}\\
------------------------------------------------------------\\
The agent acts as a message broker between the raw socket inspector and the encrypted remote network backend. It implements hot-reloading configurations and real-time alert processing.

\textbf{4.1 Message Severity and Routing Policy}\\
The agent processes alerts based on severity profiles:

\begin{center}
\begin{tabular}{L{3.5cm} C{1.5cm} L{4cm} L{5.4cm}}
\hline
\thfmt{Severity} & \thfmt{Value} & \thfmt{Examples} & \thfmt{Routing} \\
\hline
CRITICAL & 2 & PHP Shell, SYN Flood       & Sent immediately (bypasses buffer) \\
WARNING  & 1 & XSS, Slowloris, Blocklist  & Sent immediately (bypasses buffer) \\
INFO     & 0 & Normal traffic metadata    & Buffered; flushed every 5 minutes \\
\hline
\end{tabular}
\end{center}

\textbf{4.2 Heartbeat Protocol and Reload Sequence}\\
A keep-alive ping (\texttt{MSG\_TYPE\_PING}) is sent when the channel is idle for 5 minutes. On receiving a \texttt{SIGHUP} signal, the agent re-reads the configuration file and flushes dynamic entries to the kernel hook:

\begin{codeblock}
# /etc/fw_inspect/blocklist.txt
# Directives and CIDR entries (one per line)
block_local_ips=0   # 0 = Allow local IPs but log threats; 1 = Block local IPs
10.0.0.0/8
192.168.1.50/32
\end{codeblock}

\vspace{1.5em}
\textbf{5. COMPONENT: SOC RECEIVER (\texttt{soc\_receiver})}\\
------------------------------------------------------------\\
The endpoint engine is written in C as a single-threaded server leveraging \texttt{epoll(7)}. It processes agent streams, decryption, database persistence, and websocket/SSE events on distinct ports.

\textbf{5.1 Decryption Pipeline}\\
Encrypted payloads received from verified agents are processed via a system execution fork using the GnuPG backend wrapper:

\begin{codeblock}
gpg --decrypt --batch --yes --output /tmp/event.bin /tmp/event.pgp
\end{codeblock}

Decrypted telemetry is extracted into memory, verified against buffer limits, and inserted into the SQLite database. Registered web browsers connected to \texttt{/api/events/live} are simultaneously notified using Server-Sent Events (SSE).

\vspace{1.5em}
\textbf{6. COMPONENT: WEB DASHBOARD}\\
------------------------------------------------------------\\
The front-end user interface is structured as a single-page application (SPA) built using native web components and vanilla CSS. It offers charts, tables, and system logs.

\textbf{6.1 UI Management Panels}\\
The console layout contains five essential operational sectors:

\begin{center}
\begin{tabular}{L{4cm} L{10.5cm}}
\hline
\thfmt{Panel} & \thfmt{Description} \\
\hline
Threat Chart       & Bar chart: event counts per threat type over time \\
Active Agents      & Lists all nodes with IPs and last-seen timestamps \\
Live Activity Feed & Real-time scrolling event feed via SSE \\
Historical Events  & Paginated event table from SQLite via /api/events \\
Stats Overview     & Total events, unique IPs, threat breakdown cards \\
\hline
\end{tabular}
\end{center}

\textbf{6.2 Application REST API Specifications}\\
Inter-system API communication utilizes the following endpoints:

\begin{center}
\begin{tabular}{C{2.0cm} L{4.2cm} L{8.3cm}}
\hline
\thfmt{Method} & \thfmt{Endpoint} & \thfmt{Description} \\
\hline
POST & /auth            & Login; returns Set-Cookie: session\_token \\
GET  & /api/\allowbreak stats       & Event counts grouped by threat\_type \\
GET  & /api/\allowbreak agents      & List of registered agent nodes \\
GET  & /api/\allowbreak events      & Historical events (limit/offset) \\
GET  & /api/\allowbreak events/\allowbreak live & SSE stream of live events \\
\hline
\end{tabular}
\end{center}

\vspace{1.5em}
\textbf{7. PGP ENCRYPTED TRANSPORT PROTOCOL}\\
------------------------------------------------------------\\
The server-agent network link is encrypted end-to-end. Public keys are generated on startup and synced using a common Docker storage volume.

\textbf{7.1 Trust Key Exchange Sequence}\\
The automated handshaking routine is defined below:

\begin{center}
\begin{tabular}{C{1.0cm} L{3.5cm} L{10.0cm}}
\hline
\thfmt{Step} & \thfmt{Container} & \thfmt{Action} \\
\hline
1 & soc-server    & Generate RSA-2048 keypair for soc@soc.local \\
2 & soc-server    & Export public key to /etc/keys/soc\_pubkey.asc \\
3 & apache-server & Generate RSA-2048 keypair for agent@soc.local \\
4 & apache-server & Import SOC public key from shared volume \\
5 & apache-server & Export agent public key to shared volume \\
6 & soc-server    & Key watcher detects agent\_pubkey.asc (SHA-256 change) and imports \\
7 & Both          & Trust established. Agent encrypts all messages with SOC key. \\
\hline
\end{tabular}
\end{center}

\textbf{7.2 Frame Structure}\\
The low-level network header prepended to the PGP ciphertext matches the following protocol struct:

\begin{codeblock}
struct soc_msg_hdr {
    uint32_t magic;           /* 0xDEADBEEF - validates message framing     */
    uint8_t  msg_type;        /* 0=PING  1=BATCH  2=ALERT                   */
    uint8_t  agent_uuid[16];  /* Unique per-agent identifier                */
    uint32_t payload_len;     /* Length of the following PGP ciphertext     */
};
\end{codeblock}

\vspace{1.5em}
\textbf{8. THREAT DETECTION SIGNATURES}\\
------------------------------------------------------------\\
Payload inspection operates on exact pattern matching targeting the HTTP parameters.

\textbf{8.1 Remote Code Execution (RCE) Signatures}\\
The stack flags known dangerous PHP executions in request paths and body arguments:

\begin{center}
\begin{tabular}{L{3.5cm} C{3cm} L{7.9cm}}
\hline
\thfmt{Signature} & \thfmt{Severity} & \thfmt{Example Attack} \\
\hline
eval(           & CRITICAL & eval(base64\_decode(\$\_POST[cmd])) \\
system(         & CRITICAL & GET /?cmd=system(id) \\
exec(           & CRITICAL & exec('cat /etc/passwd') \\
shell\_exec(    & CRITICAL & shell\_exec('rm -rf /') \\
passthru(       & CRITICAL & passthru(\$\_GET['x']) \\
popen(          & CRITICAL & popen('ls', 'r') \\
proc\_open(     & CRITICAL & proc\_open('whoami', ...) \\
base64\_decode( & CRITICAL & Obfuscated payload delivery \\
\hline
\end{tabular}
\end{center}

\textbf{8.2 Cross-Site Scripting (XSS) Patterns}\\
The following patterns are scanned case-insensitively using custom string parsers:

\begin{center}
\begin{tabular}{L{4cm} L{10.5cm}}
\hline
\thfmt{Signature} & \thfmt{Attack Vector} \\
\hline
<script>     & Direct script injection \\
javascript:  & Anchor href and meta refresh attacks \\
onerror=     & Image/object error event handler injection \\
onload=      & Body/iframe onload event injection \\
<img         & Image tag with onerror payload \\
alert(       & Proof-of-concept XSS marker \\
\hline
\end{tabular}
\end{center}

\vspace{1.5em}
\textbf{9. PROTOCOL ATTACK MITIGATION}\\
------------------------------------------------------------\\
Layer 4 attacks are tracked and blocked based on stateful monitoring counters.

\textbf{9.1 SYN Flood Thresholds}\\
Mitigation uses sliding-window tracking parameters:

\begin{center}
\begin{tabular}{L{5cm} L{2.5cm} L{7cm}}
\hline
\thfmt{Parameter} & \thfmt{Value} & \thfmt{Description} \\
\hline
SYN Flood Threshold & 100/s   & Max SYN/s per source IP before blocking \\
Window Duration     & 1 sec   & Rate counter reset interval \\
Block Duration      & 60s     & Packets dropped until pool eviction \\
Alert throttle      & 1/window & One SOC alert when threshold is crossed \\
\hline
\end{tabular}
\end{center}

\textbf{9.2 Slowloris Constraints}\\
Slow HTTP header streams are detected using connection timers:

\begin{center}
\begin{tabular}{L{5cm} L{2.5cm} L{7cm}}
\hline
\thfmt{Parameter} & \thfmt{Value} & \thfmt{Description} \\
\hline
SLOWLORIS\_TIMEOUT & 30s & Max time for HTTP headers to complete \\
Detection basis    & Timer & Set on first packet, checked on all subsequent \\
Action             & DROP  & Connection dropped; WARNING sent instantly to SOC \\
\hline
\end{tabular}
\end{center}

\vspace{1.5em}
\textbf{10. IP BLOCKLIST HASH TABLE}\\
------------------------------------------------------------\\
The blocklist table utilizes an 8-bit hash index with 256 dynamic buckets to store blacklisted hosts and CIDR ranges. The file also supports a configuration directive to define how LAN and loopback IP traffic is handled during detection events:

\begin{codeblock}
# /etc/fw_inspect/blocklist.txt
block_local_ips=0   # 0 = Bypass drop verdicts on local IPs but log events
10.0.0.0/8
192.168.1.50/32
45.33.32.0/24
\end{codeblock}

If \texttt{block\_local\_ips} is set to 0 (default), threats originating from local private networks (RFC 1918: 10.0.0.0/8, 172.16.0.0/12, 192.168.0.0/16, 127.0.0.0/8) are allowed to pass through, but the alert is still sent to the SOC server for analysis. If set to 1, local IP threats are dropped. SIGHUP triggers live updates to this configuration without restarts.

\vspace{1.5em}
\textbf{11. AUTHENTICATION \& ACCESS CONTROL}\\
------------------------------------------------------------\\
Access rules are enforced before requests hit the application layer.

\textbf{11.1 Allowed Administrator IP Ranges}\\
Only the following network origins can access port 8443:

\begin{center}
\begin{tabular}{L{4cm} L{10.5cm}}
\hline
\thfmt{Range} & \thfmt{Purpose} \\
\hline
127.0.0.1/32   & Localhost administration \\
172.16.0.0/12  & Docker bridge and internal ranges \\
192.168.0.0/16 & LAN subnet for SOC analysts \\
\hline
\end{tabular}
\end{center}

\textbf{11.2 Credentials Configuration}\\
Administrators store hashed passwords in a localized flat file. Successful login returns an HMAC signature token:

\begin{codeblock}
# /etc/soc/users.conf format
admin:$2b$12$<bcrypt_hash>

# Session cookie format
session_token=username:unix_timestamp:hmac_sha256_hex
\end{codeblock}

\vspace{1.5em}
\textbf{12. DOCKER DEPLOYMENT SPECS}\\
------------------------------------------------------------\\
Container constraints are mapped using the compose configuration schema:

\begin{center}
\begin{tabular}{L{3.5cm} L{5.5cm} L{5.5cm}}
\hline
\thfmt{Parameter} & \thfmt{apache-server} & \thfmt{soc-server} \\
\hline
Base Image    & debian:bullseye-slim & debian:bullseye-slim \\
Exposed Ports & 80 (HTTP)            & 1113 (agent), 8443 (web) \\
Capabilities  & NET\_ADMIN (iptables) & None \\
Volumes       & ./signatures, ./keys & ./keys, soc-data (SQLite) \\
Restart       & unless-stopped       & unless-stopped \\
\hline
\end{tabular}
\end{center}

\textbf{12.1 Manual Operational Triggers}\\
Deployment and control actions are initiated using the commands below:

\begin{codeblock}
# Full clean start (wipes database and keys)
docker compose down -v
docker compose up -d --build

# Hot-reload blocklist without restart
docker exec apache-2.4.26 kill -HUP $(pidof fw_agent)
\end{codeblock}

\vspace{1.5em}
\textbf{13. TESTING AND ENGINE VALIDATION}\\
------------------------------------------------------------\\
Local validation runs a mock environment using \texttt{userspace\_compat.h} shims:

\begin{center}
\begin{tabular}{L{2.8cm} L{5.4cm} L{3.8cm} C{2.4cm}}
\hline
\thfmt{Test} & \thfmt{Input} & \thfmt{Expected} & \thfmt{Result} \\
\hline
HTTP T1      & GET /index.html              & ALLOW            & PASS \\
HTTP T2      & POST body: eval(data)         & DROP - PHP Shell & PASS \\
HTTP T3      & GET /?q=<script>alert(1)      & DROP - XSS       & PASS \\
HTTP T4      & GET /?name=\allowbreak systematic\_\allowbreak eval   & ALLOW (no FP)    & PASS \\
Blocklist T1 & 192.168.1.50 (exact)          & DROP             & PASS \\
Blocklist T2 & 192.168.1.51 (unlisted)       & ALLOW            & PASS \\
Blocklist T3 & 10.250.4.5 in 10.0.0.0/8      & DROP             & PASS \\
Blocklist T4 & 11.0.0.1 outside subnet       & ALLOW            & PASS \\
SYN T1       & 101 SYN/s from one IP         & DROP + ALERT     & PASS \\
SYN T2       & 50 SYN/s from one IP          & ALLOW            & PASS \\
Slowloris T1 & Headers incomplete after 30s  & DROP + ALERT     & PASS \\
\hline
\end{tabular}
\end{center}

\vspace{1.5em}
\textbf{14. SYSTEM OPERATIONS GUIDE}\\
------------------------------------------------------------\\
Daily configuration rules and references are documented below.

\textbf{14.1 Login Details}\\
Connection targets for security analysts:

\begin{center}
\begin{tabular}{L{4cm} L{10.5cm}}
\hline
\thfmt{Field} & \thfmt{Value} \\
\hline
Local URL  & http://127.0.0.1:8443/ \\
LAN URL    & http://<your-lan-ip>:8443/ \\
Username   & admin \\
Password   & password \\
\hline
\end{tabular}
\end{center}

\newpage
\textbf{14.2 Threat Types Reference}\\
Telemetry codes returned by the threat detection system:

\begin{center}
\begin{tabular}{C{2.5cm} L{3cm} C{3.5cm} L{5cm}}
\hline
\thfmt{threat\_\allowbreak type} & \thfmt{Name} & \thfmt{Severity} & \thfmt{Description} \\
\hline
1 & PHP Shell    & CRITICAL & RCE via PHP built-in functions \\
2 & XSS          & WARNING  & Cross-site scripting payload detected \\
3 & SYN Flood    & CRITICAL & SYN packet rate exceeds threshold \\
4 & Slowloris    & WARNING  & Incomplete HTTP header connection timeout \\
5 & IP Blocklist & WARNING  & Connection from known malicious IP \\
\hline
\end{tabular}
\end{center}

\textbf{14.3 Maintenance Command Cheat-Sheet}\\
Operational utilities for log monitoring and manual configuration:

\begin{codeblock}
# Firewall verdicts and alerts
docker exec apache-2.4.26 cat /var/log/fw_nfq.log

# Agent routing and SOC transmissions
docker exec apache-2.4.26 cat /var/log/fw_agent.log

# SOC receiver log
docker logs soc-dashboard

# Query events via API
curl -b cookies.txt http://127.0.0.1:8443/api/events

# Add IPs to blocklist then hot-reload
notepad signatures/blocklist.txt
docker exec apache-2.4.26 kill -HUP $(pidof fw_agent)
\end{codeblock}

============================================================\\
\textbf{END OF REPORT}\\
============================================================

\end{document}
